\section{State of the Art}
Traditional conceptions of mobility were primarily based on metrics that fail to capture all the factors at play. The work of Geurs and van Wee \cite{geurs_accessibility_2004} provides a comprehensive review of accessibility measures. They identified a gap in operational planning methods, which tend to be overly reductionist by using metrics such as travel speed or congestion level to facilitate the interpretation and communication of results. They showed that the quality and theoretical soundness of such evaluations could be greatly improved by taking into account behavioral and spatial dimensions, as well as the feedback loops that shape mobility, but their framework still does not account for the lived experience of travel.\\

This is where sociological approaches, from Sheller and Urry \cite{sheller_new_2006} and more notably Kaufmann's concept of motility \cite{kaufmann_motility_2004} and Doherty's notion of viscosity \cite{doherty_agentive_2015}, shine by providing a theoretical framework to explain the mismatch between predicted theoretical outcomes and observed results. Their work emphasizes how individuals experience travel and how their choices are shaped by these experiences. Jain and Lyons \cite{jain_gift_2008} further contribute to this through their research “The Gift of Travel Time,” by challenging the idea that travel time is purely a waste and should therefore be minimized. They demonstrate how individuals can re-appropriate this time for meaningful activities, making perceived time fundamentally different from measured time. This distinction shows that the quality of the journey matters at least as much as its speed. Bühler \cite{buhler_habit_2020} extends this perspective by noting the importance of habits and routines in modal choice, illustrating how decision-making is not purely rational but embedded in social and temporal practices.\\

Together, these works depict a more complete image of mobility through, first, the classical approach, which models accessibility as a measurable property, and second, the sociological approach, which presents mobility as a lived and experienced process.\\

The present study aims to assess how far these theoretical concepts have made their way into operational planning by analyzing a corpus of Swiss spatial policy documents, in order to determine whether planning discourse reflects a human-centered understanding of mobility or remains anchored in purely functional and physical-based reasoning. To do so, a total of fourty word where selected from the vocabulary present in theses publication and through exploration of the corpus and then mapped to a concept. With a few words per concept we can start to measure how much a concept is present in a text corpus as theorized in the textometric approach.\\